\documentclass[border=0pt]{standalone}
% The shared preamble of every docs diagram (docs/diagrams/<name>.tex).
% dvisvgm's driver first, so patterns and opacity survive into the SVG.
\def\pgfsysdriver{pgfsys-dvisvgm.def}
\usepackage[siunitx,RPvoltages]{circuitikz}
\usepackage{xcolor}
% Latin Modern: outline fonts for every glyph a label uses, µ included.
\usepackage[T1]{fontenc}
\usepackage{lmodern}
\usetikzlibrary{arrows.meta,backgrounds,calc,patterns,positioning}
\tikzset{
  % A white card, so diagrams read in mdBook's dark themes too.
  every picture/.append style={show background rectangle,
    background rectangle/.style={fill=white,inner sep=4pt}},
  % Layout geometry: layers told apart by fill and hatching, never colour alone.
  layer a/.style={draw=black,line width=0.6pt,fill=gray!20},
  layer b/.style={draw=black,line width=0.6pt,fill=gray!45},
  layer c/.style={draw=black,line width=0.6pt,fill=white,postaction={pattern=north east lines}},
  layer d/.style={draw=black,line width=0.6pt,fill=white,postaction={pattern=crosshatch dots}},
  marker/.style={draw=black,dashed,line width=0.6pt},
  % A measured distance and its label.
  dim/.style={{Latex[length=2mm]}-{Latex[length=2mm]},line width=0.5pt},
  dim label/.style={font=\small,fill=white,inner sep=1pt},
  % Where a rule fails: thick, dashed and labelled, not only red.
  violation/.style={draw=red!75!black,line width=1.4pt,dash pattern=on 3pt off 2pt},
  violation label/.style={font=\small\bfseries,text=red!75!black},
  caption/.style={font=\small},
}

\begin{document}
\begin{circuitikz}
  % An inverter whose output feeds its own input.
  \node[pmos] (P) at (2,2) {};
  \node[nmos] (N) at (2,0) {};
  \draw (P.S) -- ++(0,0.3) node[vcc] {VDD};
  \draw (N.S) -- ++(0,-0.3) node[ground] {};
  \node[right] at ($(N.S)+(0.15,-0.55)$) {VSS};
  \draw (P.D) -- (N.D);
  \draw (P.G) -- ++(-0.4,0) coordinate (gp) -- (gp |- N.G) -- (N.G);
  \draw (P.D |- 0,1) node[circ] {} -- ++(0.8,0) coordinate (y);
  \draw[violation] (y) -- ++(0,-3.0) -| ($(gp |- 0,1)+(-0.6,0)$) -- (gp |- 0,1);
  \draw (gp |- 0,1) node[circ] {};
  \node[right] at (y) {Y};
  \node[violation label, below] at ($(y)+(-1.2,-3.0)$) {Y drives its own gate};
\end{circuitikz}
\end{document}
